\section*{अध्याय \ref{ch:b2:deep-time-climate} --- भूवैज्ञानिक समय में जलवायु}

\begin{solution}{exo:b2:deep-time-climate:1}
(क) $900 + 37100 + 700 + 3 = \qty{38703}{PgC}$, वायु के \qty{591}{PgC} का लगभग 65 गुना।
(ख) $\num{3.5e6}\times0.012011 = \qty{42000}{PgC}$। (ग) $591/42000 = \qty{1.4}{\percent}$।
\end{solution}

\begin{solution}{exo:b2:deep-time-climate:2}
(क) निवेश $11.3 + 3 + 12.3 + 8.1 + 2.0 = \qty{36.7}{Tmol.yr^{-1}}$; निर्गम $26 + 16 =
\qty{42}{Tmol.yr^{-1}}$, अपनी अनिश्चितताओं के भीतर बराबर। (ख) $(11.3 + 3)/36.7 = 0.39$: आधे से
काफ़ी कम; निवेश का अधिकांश भाग पुराना अवसादी कार्बन है। (ग) $\num{3.5e6}/42 = \qty{8.3e4}{yr}$।
\end{solution}

\begin{solution}{exo:b2:deep-time-climate:3}
(क) \ce{Mg2SiO4 + 4 CO2 + 2 H2O -> 2 Mg^{2+} + 4 HCO3^- + SiO2}। (ख) \ce{2 Mg^{2+} + 4 HCO3^-
-> 2 MgCO3 + 2 CO2 + 2 H2O} के साथ, शुद्ध अभिक्रिया \ce{Mg2SiO4 + 2 CO2 -> 2 MgCO3 + SiO2} है: प्रति मोल
फ़ॉर्स्टेराइट दो मोल कार्बन डाइऑक्साइड। (ग) $\qty{1}{PgC} = 10^{15}/12.011 =
\qty{8.33e13}{mol}$ कार्बन, जिसके लिए \qty{4.16e13}{mol} फ़ॉर्स्टेराइट चाहिए, $4.16\times10^{13}
\times\qty{140.7}{g} = \qty{5.86}{Pg}$, लगभग छह अरब टन।
\end{solution}

\begin{solution}{exo:b2:deep-time-climate:4}
हिर्नेन्शियन हिमनदन, 445 से \qty{443}{Ma}, हिमगृह; उत्तर पैलियोज़ोइक का सबसे बड़ा हिम, 307 से
\qty{299}{Ma}, हिमगृह; ट्यूरोनियन, 94 से \qty{89}{Ma}, ग्रीनहाउस (फ़ैनेरोज़ोइक का सबसे उष्ण);
PETM, 55.9 से \qty{55.7}{Ma}, ग्रीनहाउस; EECO, 53 से \qty{49}{Ma}, ग्रीनहाउस; अंटार्कटिक हिमचादर
का आरंभ, लगभग \qty{34}{Ma}, वर्तमान हिमगृह की शुरुआत; उत्तरी हिमचादरें,
\qty{2.7}{Ma}, हिमगृह।
\end{solution}

\begin{solution}{exo:b2:deep-time-climate:5}
(क) $63000/(8.314\times288^2) = \qty{0.0914}{K^{-1}}$, जो 0.09 तक पूर्णांकित होता है। (ख) $74000/(8.314
\times288^2) = \qty{0.107}{K^{-1}}$। (ग) $\ln1.2/0.09 = \qty{2.03}{K}$ और $\ln 1.2/0.107 =
\qty{1.70}{K}$। (घ) स्थिर अवस्था में अपक्षय को बहिर्गैसन के बराबर होना ही है, जो $e^{k
\Delta T}$ को और इस प्रकार $\Delta T$ को निश्चित करता है; संवेदनशीलता केवल यह तय करती है कि कितनी
कार्बन डाइऑक्साइड वह तापन उत्पन्न करती है।
\end{solution}

\begin{solution}{exo:b2:deep-time-climate:6}
यहाँ $k_T = 0.0914$। (क) सटीक $\exp[7577.5\times(1/288 - 1/293)] = 1.567$; सन्निकटन
$e^{0.0914\times5} = 1.579$: \qty{0.8}{\percent} अधिक। (ख) सटीक $e^{2.482} = 11.97$;
सन्निकटन $e^{2.741} = 15.50$: \qty{30}{\percent} अधिक। (ग) $\ln(\text{अनुपात}) = (E_{\text
a}/R)[x/T_0^2 - x/(T_0(T_0 + x))] = (E_{\text a}/R)\,x^2/(T_0^2(T_0 + x))$। इसे $\ln1.1$ के बराबर रखने पर:
$x^2 = 1.043\,(288 + x)$, अतः $x = \qty{17.9}{K}$। लगभग \qty{18}{K} से अधिक तापन पर
घातांकी रूप अपक्षय को दसवें भाग से अधिक बढ़ाकर बताता है।
\end{solution}

\begin{solution}{exo:b2:deep-time-climate:7}
(क) $1444/276 = 5.23$, और $3.93\log_2 5.23 = \qty{9.4}{W.m^{-2}}$। (ख) लगभग ठीक सूर्य के
\qty{9}{W.m^{-2}} जितना: दोनों प्रणोदन एक-दूसरे को निरस्त करते हैं। (ग) $(1444 - 276)/419.5 =
\qty{2.8}{ppm}$ प्रति मिलियन वर्ष। \qty{3.4}{ppm} सभी आँकड़ों पर समंजित एक सीधी रेखा से आता है,
जो समतलित वक्र के दो छोरों के बजाय हर अनुमान को भार देती है।
\end{solution}

\begin{solution}{exo:b2:deep-time-climate:8}
(क) $3000/20000 = \qty{0.15}{PgC.yr^{-1}}$ से $7000/3000 = \qty{2.3}{PgC.yr^{-1}}$ तक, एक परास
जिसमें IPCC के 0.3 से 1.5 समाहित हैं। (ख) $35\times12.011/44.01 = \qty{9.6}{PgC.yr^{-1}}$, PETM के
0.3 से \qty{1.5}{PgC.yr^{-1}} का 6 से 32 गुना। (ग) $3.93\log_2(2000/900) = \qty{4.5}{W.m^{-2}}$,
और $3\log_2 2.22 = \qty{3.5}{\celsius}$: \qty{5}{\celsius} का लगभग सत्तर प्रतिशत। अन्य
प्रतिपुष्टियों या तापन के किसी अन्य स्रोत, जैसे मीथेन, को शेष की पूर्ति करनी होगी।
\end{solution}

\begin{solution}{exo:b2:deep-time-climate:9}
(क) अब $1.2\,W_0e^{k\Delta T} = V = W_0$, अतः $\Delta T = -\ln 1.2/0.09 = \qty{-2.03}{K}$।
(ख) $278\times2^{-2.03/3} = \qty{174}{ppm}$। (ग) नहीं: किसी भी स्थिर अवस्था में अपक्षय बहिर्गैसन के बराबर
होता है, जो बदला नहीं है। ऊँची अपक्षयशीलता की भरपाई अधिक ठंडी जलवायु करती है।
\end{solution}

\begin{solution}{exo:b2:deep-time-climate:10}
(क) $k = 0.119$: $\tau = 591\times0.693/(0.08\times0.119\times3\times0.138) = \qty{104}{kyr}$।
(ख) $S = 8$: $\tau = \qty{51.5}{kyr}$। (ग) दोनों तापस्थापी को तेज़ बनाते हैं, क्योंकि $\tau \propto
1/(kS)$। $E_{\text a}$ लगभग दो के गुणक से अनिश्चित है, और इन समय-पैमानों पर $S$ भी लगभग उतने ही
गुणक से (3 बनाम 5 से \qty{8}{K}): दोनों लगभग बराबर महत्त्व रखते हैं।
\end{solution}

\begin{solution}{exo:b2:deep-time-climate:11}
(क) $0.08\times5000 = \qty{400}{PgC}$ वायु में रहता है: $C/C_0 = 991/591 = 1.677$, यानी
\qty{466}{ppm}; $\Delta T = 3\log_2 1.677 = \qty{2.24}{K}$। (ख) $W - W_0 = 0.138\,(e^{0.09\times2.24}
- 1) = \qty{0.0308}{PgC.yr^{-1}}$, \qty{22}{\percent} अधिक अपक्षय; $5000/0.0308 =
\qty{162}{kyr}$। (ग) $\tau = \qty{137}{kyr}$ और $\tau\ln10 = \qty{315}{kyr}$। (घ) तापन कार्बन
डाइऑक्साइड के केवल लघुगणक के रूप में बढ़ता है, अतः एक बड़ा अतिरिक्त अपक्षय को रैखिक नियम की
मान्यता से कम बढ़ाता है; निष्कासन अधिक धीमे शुरू होता है।
\end{solution}

\begin{solution}{exo:b2:deep-time-climate:12}
(क) $M^*$ के निकट $\mathrm d\Delta T/\mathrm dM = Sf/(A_0\ln 2\,(1 + fM^*/A_0)) = Sf/(A_0\ln 2\,
C^*/C_0)$ और $\mathrm dW/\mathrm dM = kW\,\mathrm d\Delta T/\mathrm dM$, जहाँ वहाँ $W = V$ है;
इसका व्युत्क्रम $\tau'$ है। (ख) $C^*/C_0 = 2^{2.03/3} = 1.597$: $\tau' = 137.4\times1.597/1.2 =
\qty{183}{kyr}$। (ग) धीमा: ऊँची कार्बन डाइऑक्साइड पर जोड़ा गया हर पेटाग्राम कार्बन कम गर्म करता है,
क्योंकि प्रणोदन लघुगणकीय है, अतः अपक्षय उसके प्रति कम अनुक्रिया देता है।
\end{solution}

\begin{solution}{pb:b2:deep-time-climate:1}
\textbf{1.} $1.3 + 1.3 + 4.5 + 4.2 = \qty{11.3}{Tmol.yr^{-1}}$, $\qty{0.136}{PgC.yr^{-1}}$।
\textbf{2.} $11.5\times0.012011 = \qty{0.138}{PgC.yr^{-1}}$।
\textbf{3.} वे लगभग \qty{2}{\percent} के भीतर मेल खाते हैं: सिलिकेट अपक्षय ज्वालामुखीय बहिर्गैसन को
संतुलित करता है, जैसा \omterm{def:b2:deep-time-climate:carbon-cycle}{दीर्घकालिक कार्बन चक्र} अपेक्षित करता है।
\textbf{4.} $\num{5.1e18}/0.02897 = \qty{1.76e20}{mol}$ वायु; $\num{591e15}/12.011 =
\qty{4.92e16}{mol}$ कार्बन; $\num{4.92e16}/\num{1.76e20} = \num{2.80e-4}$, \qty{280}{ppm}।
\textbf{5.} $591/0.136 = \qty{4350}{yr}$।
\textbf{6.} $1 + 0.08 = 1.08$: केवल \qty{8}{\percent} अधिक।
\textbf{7.} $63000/(8.314\times288^2) = \qty{0.0914}{K^{-1}}$, जो 0.09 तक पूर्णांकित होता है।
\textbf{8.} $e^{0.09} = 1.09$: प्रति केल्विन \qty{9}{\percent} अधिक।
\textbf{9.} $\Delta T = S\log_2(1 + fM/A_0)$; 0 पर इसका अवकलज $Sf/(A_0\ln2) =
3\times0.08/(591\times0.693) = \qty{5.86e-4}{K.PgC^{-1}}$ है।
\textbf{10.} $W_0k\times\num{5.86e-4} = 0.138\times0.09\times\num{5.86e-4} = \qty{7.28e-6}{yr^{-1}}$।
\textbf{11.} इसका व्युत्क्रम $A_0\ln2/(fkSW_0)$ है, प्रतिज्ञप्ति का $\tau$, क्योंकि
रैखिकीकृत समीकरण $\mathrm dM/\mathrm dt = -(\mathrm dW/\mathrm dM)\,M$ है।
\textbf{12.} चारों: $\tau$ उनमें से प्रत्येक के व्युत्क्रमानुपाती है।
\textbf{13.} $1/\num{7.28e-6} = \qty{1.37e5}{yr}$।
\textbf{14.} $137\times0.09/0.119 = \qty{104}{kyr}$।
\textbf{15.} $k_T = 0.107$: $137\times0.09/0.107 = \qty{115}{kyr}$।
\textbf{16.} 1.5 से 7 के गुणक से छोटे: एक-बॉक्स मॉडल और उसका स्थिर $f$ तापस्थापी को
बहुत तेज़ बना देते हैं।
\textbf{17.} $\num{50e6}/\num{1.37e5} = 364$: तापस्थापी के पास हर चरण पर चक्र को संतुलित करने का
समय था, अतः गिरावट बहिर्गैसन या अपक्षयशीलता के धीमे परिवर्तनों के पीछे चली।
\textbf{18.} $16\times0.137 = \qty{2.2}{ppm}$।
\textbf{19.} $0.08\times5000 = \qty{400}{PgC}$; $278\times991/591 = \qty{466}{ppm}$।
\textbf{20.} $3\log_2 1.677 = \qty{2.24}{K}$।
\textbf{21.} $0.138\,(e^{0.202} - 1) = \qty{0.0308}{PgC.yr^{-1}}$, \qty{22.3}{\percent}।
\textbf{22.} $137\times\ln10 = \qty{315}{kyr}$।
\textbf{23.} अपक्षय अधिक संवेदनशील रहा हो सकता है (अपवाह, बड़ी सक्रियण ऊर्जा, बड़ी पृथ्वी-तंत्र
संवेदनशीलता की धीमी प्रतिपुष्टियाँ); और स्थल तथा समुद्र का कार्बनिक कार्बन फिर से बढ़कर
सिलिकेट अपक्षय से तेज़ी से कार्बन को वायु से निकाल सकता था।
\textbf{24.} इस मॉडल में \omterm{def:b2:deep-time-climate:thermostat}{सिलिकेट-अपक्षय तापस्थापी} अतिरिक्त कार्बन को लगभग
\textbf{\qty[text-series-to-math=true]{137}{} हज़ार वर्ष} के e-गुना क्षय काल के साथ हटाता है।
\end{solution}
